\chapter{2013年四川卷（文）}

\section{单选题}

\begin{problem}
设集合 \(A = \{1, 2, 3\}\)，集合 \(B = \{-2, 2\}\)，则 \(A \cap B =\)（\quad）

\choices
    {\(\varnothing\)}
    {\(\{2\}\)}
    {\(\{-2, 2\}\)}
    {\(\{-2, 1, 2, 3\}\)}
\end{problem}

\begin{answer}
B
\end{answer}

\begin{solution}
交集中的元素必须同时属于 \(A\) 和 \(B\). 集合 \(A\) 中只有 \(2\) 也属于集合 \(B\)，所以 \(A\cap B=\{2\}\)，故选 B.
\end{solution}

\begin{problem}
一个几何体的三视图如图所示，则该几何体可以是（\quad）

\bitmapfigure[max width=0.32\linewidth]{img/2013/sichuan_liberal/q2_fig1.png}

\choices
    {棱柱}
    {棱台}
    {圆柱}
    {圆台}
\end{problem}

\begin{answer}
D
\end{answer}

\begin{solution}
正视图和侧视图都是等腰梯形，俯视图是两个同心圆组成的圆环. 因此，该几何体的上下底面是两个半径不同且同心的圆，侧面为梯形，符合圆台的三视图特征，故选 D.
\end{solution}

\begin{problem}
如图，在复平面内，点 \(A\) 表示复数 \(z\)，则图中表示 \(z\) 的共轭复数的点是（\quad）

\bitmapfigure[max width=0.28\linewidth]{img/2013/sichuan_liberal/q3_fig1.png}

\choices
    {\(A\)}
    {\(B\)}
    {\(C\)}
    {\(D\)}
\end{problem}

\begin{answer}
B
\end{answer}

\begin{solution}
复数 \(z\) 与其共轭复数 \(\overline z\) 的实部相同、虚部互为相反数，所以在复平面内，表示它们的两点关于实轴对称. 图中点 \(A\) 关于实轴的对称点是 \(B\)，故选 B.
\end{solution}

\begin{problem}
设 \(x \in \Z\)，集合 \(A\) 是奇数集，集合 \(B\) 是偶数集. 若命题 \(p: \forall x \in A\)，\(2x \in B\)，则（\quad）

\choices
    {\(\neg p: \exists x \in A, 2x \in B\)}
    {\(\neg p: \exists x \notin A\)，\(2x \in B\)}
    {\(\neg p: \exists x \in A\)，\(2x \notin B\)}
    {\(\neg p: \forall x \notin A\)，\(2x \notin B\)}
\end{problem}

\begin{answer}
C
\end{answer}

\begin{solution}
命题 \(p\) 是“对所有 \(x\in A\)，都有 \(2x\in B\)”. 全称命题的否定是存在一个反例，因此
\(\neg p: \exists x\in A\)，使得 \(2x\notin B\)，故选 C.
\end{solution}

\begin{problem}
抛物线 \(y^{2}=8x\) 的焦点到直线 \(x-\sqrt{3}y=0\) 的距离是（\quad）

\choices
    {\(2\sqrt{3}\)}
    {\(2\)}
    {\(\sqrt{3}\)}
    {\(1\)}
\end{problem}

\begin{answer}
D
\end{answer}

\begin{solution}
将抛物线写成 \(y^{2}=2px\)，得 \(2p=8\)，所以 \(p=4\)，焦点为 \(F\left(\frac p2,0\right)=(2,0)\). 点 \((x_0,y_0)\) 到直线 \(x-\sqrt{3}y=0\) 的距离为
\(\dfrac{|x_0-\sqrt{3}y_0|}{\sqrt{1+3}}\)，故焦点到该直线的距离为 \(\dfrac{|2|}{2}=1\)，故选 D.
\end{solution}

\begin{problem}
函数 \(f(x)=2\sin\left(\omega x+\varphi\right)\)（\(\omega>0\)，\(-\frac{\pi}{2}<\varphi<\frac{\pi}{2}\)）的部分图象如图所示，则 \(\omega\)，\(\varphi\) 的值分别是（\quad）

\bitmapfigure[max width=0.34\linewidth]{img/2013/sichuan_liberal/q6_fig1.png}

\choices
    {\(2, -\frac{\pi}{3}\)}
    {\(2\)，\(-\frac{\pi}{6}\)}
    {\(4\)，\(-\frac{\pi}{6}\)}
    {\(4\)，\(\frac{\pi}{3}\)}
\end{problem}

\begin{answer}
A
\end{answer}

\begin{solution}
由图象可知，相邻的最大值点和最小值点的横坐标分别为 \(\frac{5\pi}{12}\) 和 \(\frac{11\pi}{12}\)，两者相距半个周期. 因此 \(\frac{T}{2}=\frac{\pi}{2}\)，得 \(T=\pi\)，从而 \(\omega=\frac{2\pi}{T}=2\). 在最大值点 \(x=\frac{5\pi}{12}\) 处，有
\(2\cdot\frac{5\pi}{12}+\varphi=\frac\pi2+2k\pi\). 结合 \(-\frac\pi2<\varphi<\frac\pi2\)，取 \(k=0\)，得 \(\varphi=-\frac\pi3\). 所以 \(\omega=2\)，\(\varphi=-\frac\pi3\)，故选 A.
\end{solution}

\begin{problem}
某学校随机抽取 20 个班，调查各班中有网上购物经历的人数，所得数据的茎叶图如图所示，以组距为 5 将数据分组成 \([0, 5)\)，\([5, 10)\)，\(\cdots\)，\([30, 35)\)，\([35, 40]\) 时，所作的频率分布直方图是（\quad）
\begin{center}
\begin{tabular}{r|l}
\cline{2-2}
0 & \(7\quad 3\) \\
1 & \(7\)，\(6\)，\(4\)，\(4\)，\(3\)，\(0\) \\
2 & \(7\)，\(5\)，\(5\)，\(4\)，\(3\)，\(2\)，\(0\) \\
3 & \(8\)，\(5\)，\(4\)，\(3\)，\(0\)
\end{tabular}
\end{center}

\choices
    {\choicetexfigure[width=0.15\paperwidth]{img_repaint/2013/sichuan_liberal/q7_choice_a.tex}}
    {\choicetexfigure[width=0.15\paperwidth]{img_repaint/2013/sichuan_liberal/q7_choice_b.tex}}
    {\choicetexfigure[width=0.15\paperwidth]{img_repaint/2013/sichuan_liberal/q7_choice_c.tex}}
    {\choicetexfigure[width=0.15\paperwidth]{img_repaint/2013/sichuan_liberal/q7_choice_d.tex}}
\end{problem}

\begin{answer}
A
\end{answer}

\begin{solution}
按茎叶图逐组统计，落在 \([0,5)\)，\([5,10)\)，\([10,15)\)，\([15,20)\)，\([20,25)\)，\([25,30)\)，\([30,35)\)，\([35,40]\) 内的数据个数分别为 \(1\)，\(1\)，\(4\)，\(2\)，\(4\)，\(3\)，\(3\)，\(2\). 其中 \(3\) 落在 \([0,5)\) 内，\(7\) 落在 \([5,10)\) 内.

由于组距为 \(5\)，频率分布直方图的纵坐标是频率除以组距，所以各组的频率密度依次为
\(\frac{1}{20\times5}=0.01\)，\(\frac{1}{20\times5}=0.01\)，\(\frac{4}{20\times5}=0.04\)，\(\frac{2}{20\times5}=0.02\)，\(\frac{4}{20\times5}=0.04\)，\(\frac{3}{20\times5}=0.03\)，\(\frac{3}{20\times5}=0.03\)，\(\frac{2}{20\times5}=0.02\). 与选项比较，只有第一幅图符合，故选 A.
\end{solution}

\begin{problem}
若变量 \(x\)，\(y\) 满足约束条件 \(\begin{cases}
x+y\leqslant8,\\
2y-x\leqslant4,\\
x\geqslant0,\\
y\geqslant0
\end{cases}\) 且 \(z=5y-x\) 的最大值为 \(a\)，最小值为 \(b\)，则 \(a-b\) 的值是（\quad）

\choices
    {\(48\)}
    {\(30\)}
    {\(24\)}
    {\(16\)}
\end{problem}

\begin{answer}
C
\end{answer}

\begin{solution}
可行域由四个不等式确定. 其顶点为 \((0,0)\)、\((8,0)\)、\((4,4)\) 和 \((0,2)\)：其中 \((4,4)\) 是直线 \(x+y=8\) 与 \(2y-x=4\) 的交点，\((0,2)\) 是直线 \(x=0\) 与 \(2y-x=4\) 的交点.

在这些顶点处分别计算 \(z=5y-x\)，得到
\(z(0,0)=0\)，\(z(8,0)=-8\)，\(z(4,4)=16\)，\(z(0,2)=10\). 因此 \(a=16\)，\(b=-8\)，所以 \(a-b=16-(-8)=24\)，故选 C.
\end{solution}

\begin{problem}
从椭圆 \(\frac{x^2}{a^2}+\frac{y^2}{b^2}=1\)（\(a>b>0\)）上一点 \(P\) 向 \(x\) 轴作垂线，垂足恰为左焦点 \(F_1\)，\(A\) 是椭圆与 \(x\) 轴正半轴的交点，\(B\) 是椭圆与 \(y\) 轴正半轴的交点，且 \(AB \parallel OP\)（\(O\) 是坐标原点），则该椭圆的离心率是（\quad）

\choices
    {\(\frac{\sqrt{2}}{4}\)}
    {\(\frac{1}{2}\)}
    {\(\frac{\sqrt{2}}{2}\)}
    {\(\frac{\sqrt{3}}{2}\)}
\end{problem}

\begin{answer}
C
\end{answer}

\begin{solution}
设椭圆的半焦距为 \(c\)，则 \(c^2=a^2-b^2\)，左焦点为 \(F_1(-c,0)\). 因为点 \(P\) 向 \(x\) 轴的垂足是 \(F_1\)，可设 \(P=(-c,y_P)\). 将其代入椭圆方程，得
\(\frac{c^2}{a^2}+\frac{y_P^2}{b^2}=1\)，从而 \(y_P^2=\frac{b^4}{a^2}\).

又因为 \(A=(a,0)\)，\(B=(0,b)\)，所以 \(\overrightarrow{AB}=(-a,b)\). 由 \(AB\parallel OP\)，有
\((-a)y_P-b(-c)=0\)，即 \(ay_P=bc\). 因此 \(y_P>0\)，且 \(y_P=\frac{bc}{a}\). 与 \(y_P^2=\frac{b^4}{a^2}\) 比较，得 \(c=b\). 于是 \(a^2=b^2+c^2=2b^2\)，椭圆的离心率为 \(e=\frac ca=\frac{\sqrt{2}}{2}\)，故选 C.
\end{solution}

\begin{problem}
设函数 \(f(x)=\sqrt{\e^x+x-a}\)（\(a\in\R\)，\(\e\) 为自然对数的底数）. 若存在 \(b\in[0,1]\) 使 \(f(f(b))=b\) 成立，则 \(a\) 的取值范围是（\quad）

\choices
    {\([1,\e]\)}
    {\([1,1+\e]\)}
    {\([\e,1+\e]\)}
    {\([0,1]\)}
\end{problem}

\begin{answer}
A
\end{answer}

\begin{solution}
令 \(u=f(b)\)，则 \(u\geqslant0\)，并且
\(u^2=\e^b+b-a\)，\(b^2=\e^u+u-a\). 两式相减，得到
\(\e^b+b-\e^u-u=u^2-b^2\).

令 \(g(t)=\e^t+t\)，则 \(g'(t)=\e^t+1>0\)，所以 \(g(t)\) 在 \(\R\) 上严格递增. 如果 \(u>b\)，左边 \(g(b)-g(u)<0\)，右边 \(u^2-b^2>0\)，矛盾；如果 \(u<b\)，左边为正而右边为负，也矛盾. 因此只能有 \(u=b\). 代回得
\(a=\e^b+b-b^2\).

令 \(h(b)=\e^b+b-b^2\). 在 \([0,1]\) 上，利用 \(\e^b\geqslant1+b\)，有 \(h'(b)=\e^b+1-2b\geqslant2-b>0\)，所以 \(h\) 在 \([0,1]\) 上递增. 于是 \(h(0)=1\)，\(h(1)=\e\)，故 \(a\in[1,\e]\)，选 A.

反过来，任意 \(a\in[1,\e]\) 都可由函数 \(h\) 在 \([0,1]\) 上的连续性和单调性取到，即存在 \(b\in[0,1]\) 使 \(a=h(b)\). 此时 \(\e^b+b-a=b^2\)，所以 \(f(b)=\sqrt{b^2}=b\)，进而 \(f(f(b))=b\). 因此 \(a\) 的取值范围确为 \([1,\e]\)，选 A.
\end{solution}

\section{填空题}

\begin{problem}
\(\lg\sqrt{5}+\lg\sqrt{20}\) 的值是\fillinblank{}.
\end{problem}

\begin{answer}
\(1\)
\end{answer}

\begin{solution}
由对数的运算性质，\(\lg\sqrt{5}+\lg\sqrt{20}=\lg\left(\sqrt{5}\sqrt{20}\right)=\lg 10=1\)，所以填 \(1\).
\end{solution}

\begin{problem}
在平行四边形 \(ABCD\) 中，对角线 \(AC\) 与 \(BD\) 交于点 \(O\)，\(\overrightarrow{AB}+\overrightarrow{AD}=\lambda\overrightarrow{AO}\)，则 \(\lambda=\)\fillinblank{}.

\bitmapfigure[max width=0.44\linewidth]{img/2013/sichuan_liberal/q12_fig1.png}
\end{problem}

\begin{answer}
\(2\)
\end{answer}

\begin{solution}
平行四边形的对角线互相平分，所以 \(O\) 是 \(AC\) 的中点，且 \(\overrightarrow{AO}=\frac12\overrightarrow{AC}\). 又由平行四边形法则，\(\overrightarrow{AB}+\overrightarrow{AD}=\overrightarrow{AC}\). 因此
\(\overrightarrow{AB}+\overrightarrow{AD}=2\overrightarrow{AO}\)，从而 \(\lambda=2\).
\end{solution}

\begin{problem}
已知函数 \(f(x)=4x+\frac{a}{x}\)（\(x>0\)，\(a>0\)）在 \(x=3\) 时取得最小值，则 \(a=\)\fillinblank{}.
\end{problem}

\begin{answer}
\(36\)
\end{answer}

\begin{solution}
因为 \(x>0\)，\(a>0\)，由基本不等式

\[
4x+\frac{a}{x}\geqslant 2\sqrt{4x\cdot\frac{a}{x}}=4\sqrt a
\]

等号成立当且仅当 \(4x=\frac{a}{x}\)，即 \(x=\frac{\sqrt a}{2}\). 题设说明最小值在 \(x=3\) 处取得，所以

\[
\frac{\sqrt a}{2}=3
\]

从而 \(\sqrt a=6\)，\(a=36\). 反过来，当 \(a=36\) 时，\(f(x)\geqslant24=f(3)\)，故 \(x=3\) 确实取得最小值. 因此填 \(36\).
\end{solution}

\begin{problem}
设 \(\sin 2\alpha=-\sin\alpha\)，\(\alpha\in\left(\frac{\pi}{2},\pi\right)\)，则 \(\tan 2\alpha\) 的值是\fillinblank{}.
\end{problem}

\begin{answer}
\(\sqrt{3}\)
\end{answer}

\begin{solution}
由 \(\sin2\alpha=2\sin\alpha\cos\alpha\)，原等式化为 \(2\sin\alpha\cos\alpha=-\sin\alpha\). 由于 \(\alpha\in\left(\frac\pi2,\pi\right)\)，所以 \(\sin\alpha>0\)，可约去 \(\sin\alpha\)，得 \(\cos\alpha=-\frac12\). 于是 \(\alpha=\frac{2\pi}{3}\)，从而
\(\tan2\alpha=\tan\frac{4\pi}{3}=\sqrt3\). 因此填 \(\sqrt3\).
\end{solution}

\begin{problem}
在平面直角坐标系内，到点 \(A(1,2)\)，\(B(1,5)\)，\(C(3,6)\)，\(D(7,-1)\) 的距离之和最小的点的坐标是\fillinblank{}.
\end{problem}

\begin{answer}
\((2,4)\)
\end{answer}

\begin{solution}
对任意点 \(P\)，由三角不等式有 \(PA+PC\geqslant AC\)，且 \(PB+PD\geqslant BD\)，所以
\(PA+PB+PC+PD\geqslant AC+BD\). 当且仅当 \(P\) 同时在线段 \(AC\) 和线段 \(BD\) 上时取等号，因此只需求两线段的交点.

直线 \(AC\) 的方程为 \(y=2x\)，直线 \(BD\) 的方程为 \(y=6-x\). 联立得 \(2x=6-x\)，所以 \(x=2\)，\(y=4\). 两线段确实相交于该点，故所求点的坐标为 \((2,4)\).
\end{solution}

\section{解答题}

\begin{problem}
在等比数列 \(\{a_n\}\) 中，\(a_2-a_1=2\)，且 \(2a_2\) 为 \(3a_1\) 和 \(a_3\) 的等差中项，求数列 \(\{a_n\}\) 的首项，公比及前 \(n\) 项和.
\end{problem}

\begin{solution}
设数列的首项为 \(a_1\)，公比为 \(q\)，则 \(a_2=a_1q\)，\(a_3=a_1q^2\). 由 \(a_2-a_1=2\)，得
\(a_1(q-1)=2\)，所以 \(a_1\neq0\).

因为 \(2a_2\) 是 \(3a_1\) 和 \(a_3\) 的等差中项，所以
\(2(2a_2)=3a_1+a_3\). 代入等比数列各项，得 \(4a_1q=3a_1+a_1q^2\). 由于 \(a_1\neq0\)，可约去 \(a_1\)，得到
\(q^2-4q+3=0\)，即 \((q-1)(q-3)=0\).

若 \(q=1\)，则 \(a_1(q-1)=0\)，与 \(a_1(q-1)=2\) 矛盾，所以 \(q=3\). 再由 \(a_1(q-1)=2\)，得 \(a_1=1\). 因此 \(a_n=3^{n-1}\)，前 \(n\) 项和为
\(S_n=\frac{a_1(q^n-1)}{q-1}=\frac{3^n-1}{2}\).
\end{solution}

\begin{problem}
在 \(\triangle ABC\) 中，角 \(A\)，\(B\)，\(C\) 的对边分别为 \(a\)，\(b\)，\(c\)，且 \(\cos\left(A-B\right)\cos B-\sin\left(A-B\right)\sin\left(A+C\right)=-\frac35\).
\begin{enumerate}
\item 求 \(\sin A\) 的值；
\item 若 \(a=4\sqrt2\)，\(b=5\)，求向量 \(\overrightarrow{BA}\) 在 \(\overrightarrow{BC}\) 方向上的投影.
\end{enumerate}
\end{problem}

\begin{solution}
\begin{enumerate}
\item 由于 \(A+B+C=\pi\)，所以 \(A+C=\pi-B\)，从而 \(\sin(A+C)=\sin B\). 因此题设左边为
\(\cos(A-B)\cos B-\sin(A-B)\sin B=\cos A\). 于是 \(\cos A=-\frac35\). 三角形内角满足 \(0<A<\pi\)，故 \(\sin A>0\)，所以 \(\sin A=\sqrt{1-\cos^2A}=\frac45\).

\item 由余弦定理，\(a^2=b^2+c^2-2bc\cos A\). 代入 \(a=4\sqrt2\)，\(b=5\)，\(\cos A=-\frac35\)，得
\(32=25+c^2+6c\)，即 \(c^2+6c-7=0\). 由于 \(c>0\)，所以 \(c=1\).

向量 \(\overrightarrow{BA}\) 在 \(\overrightarrow{BC}\) 方向上的投影为 \(|\overrightarrow{BA}|\cos B=c\cos B\). 再由余弦定理，\(\cos B=\frac{a^2+c^2-b^2}{2ac}=\frac{32+1-25}{8\sqrt2}=\frac{\sqrt2}{2}\). 因此所求投影为 \(1\times\frac{\sqrt2}{2}=\frac{\sqrt2}{2}\).
\end{enumerate}
\end{solution}

\begin{problem}
某算法的程序框图如图所示，其中输入的变量 \(x\) 在 \(1\)，\(2\)，\(3\)，\(\cdots\)，\(24\) 这 24 个整数中等可能随机产生.

\begin{enumerate}
\item 分别求出按程序框图正确编程运行时输出 \(y\) 的值为 \(i\) 的概率 \(P_i\)（\(i=1\)，\(2\)，\(3\)）；
\item 甲，乙两同学依据自己对程序框图的理解，各自编写程序重复运行 \(n\) 次后，统计记录了输出 \(y\) 的值为 \(i\)（\(i=1\)，\(2\)，\(3\)）的频数. 以下是甲，乙所作频数统计表的部分数据.

\begin{center}
甲的频数统计表（部分）
\renewcommand{\arraystretch}{1.2}
\begin{tabular}{c|ccc}
\hline
\begin{tabular}{c}运行次数\\\(n\)\end{tabular} & \begin{tabular}{c}输出 \(y\) 的值为\\1 的频数\end{tabular} & \begin{tabular}{c}输出 \(y\) 的值为\\2 的频数\end{tabular} & \begin{tabular}{c}输出 \(y\) 的值为\\3 的频数\end{tabular} \\
\hline
30 & 14 & 6 & 10 \\
\hline
\(\cdots\) & \(\cdots\) & \(\cdots\) & \(\cdots\) \\
\hline
2100 & 1027 & 376 & 697 \\
\hline
\end{tabular}
\end{center}

\begin{center}
乙的频数统计表（部分）
\renewcommand{\arraystretch}{1.2}
\begin{tabular}{c|ccc}
\hline
\begin{tabular}{c}运行次数\\\(n\)\end{tabular} & \begin{tabular}{c}输出 \(y\) 的值为\\1 的频数\end{tabular} & \begin{tabular}{c}输出 \(y\) 的值为\\2 的频数\end{tabular} & \begin{tabular}{c}输出 \(y\) 的值为\\3 的频数\end{tabular} \\
\hline
30 & 12 & 11 & 7 \\
\hline
\(\cdots\) & \(\cdots\) & \(\cdots\) & \(\cdots\) \\
\hline
2100 & 1051 & 696 & 353 \\
\hline
\end{tabular}
\end{center}

当 \(n=2100\) 时，根据表中的数据，分别写出甲，乙所编程序各自输出 \(y\) 的值为 \(i\)（\(i=1\)，\(2\)，\(3\)）的频率（用分数表示），并判断两位同学中哪一位所编写程序符合算法要求的可能性较大.
\end{enumerate}

\bitmapfigure[max width=0.48\linewidth]{img/2013/sichuan_liberal/q18_fig1.png}
\end{problem}

\begin{solution}
\begin{enumerate}
\item 当输入的整数 \(x\) 为奇数时，程序输出 \(y=1\). 在 \(1\) 到 \(24\) 中奇数有 \(12\) 个，所以 \(P_1=\frac{12}{24}=\frac12\).

当 \(x\) 为偶数且能被 \(3\) 整除时，程序输出 \(y=3\). 这样的数为 \(6\)，\(12\)，\(18\)，\(24\)，共有 \(4\) 个，所以 \(P_3=\frac4{24}=\frac16\). 其余偶数共有 \(12-4=8\) 个，程序输出 \(y=2\)，所以 \(P_2=\frac8{24}=\frac13\).

\item 甲所编程序在 \(n=2100\) 时的三个频率依次为 \(\frac{1027}{2100}\)，\(\frac{376}{2100}=\frac{94}{525}\)，\(\frac{697}{2100}\). 乙所编程序的三个频率依次为 \(\frac{1051}{2100}\)，\(\frac{696}{2100}=\frac{58}{175}\)，\(\frac{353}{2100}\).

按算法得到的理论概率为 \(\left(\frac12,\frac13,\frac16\right)\). 乙的三个频率与理论概率的绝对偏差分别为

\(\left|\frac{1051}{2100}-\frac12\right|=\frac1{2100}\)，\(\left|\frac{696}{2100}-\frac13\right|=\frac1{525}\)，\(\left|\frac{353}{2100}-\frac16\right|=\frac1{700}\).

甲的三个频率与理论概率的绝对偏差分别为

\(\left|\frac{1027}{2100}-\frac12\right|=\frac{23}{2100}\)，\(\left|\frac{376}{2100}-\frac13\right|=\frac{27}{175}\)，\(\left|\frac{697}{2100}-\frac16\right|=\frac{347}{2100}\).

乙的三个频率均比甲对应频率更接近理论概率，因此乙同学所编程序符合算法要求的可能性较大.
\end{enumerate}
\end{solution}

\begin{problem}
如图，在三棱柱 \(ABC-A_{1}B_{1}C_{1}\) 中，侧棱 \(AA_{1}\perp\) 底面 \(ABC\)，\(AB=AC=2AA_{1}=2\)，\(\angle BAC=120^{\circ}\)，\(D\)，\(D_{1}\) 分别是线段 \(BC\)，\(B_{1}C_{1}\) 的中点，点 \(P\) 是线段 \(AD\) 上异于端点的点.
\begin{enumerate}
\item 在平面 \(ABC\) 内，试作出过点 \(P\) 与平面 \(A_{1}BC\) 平行的直线 \(l\)，请说明理由，并证明直线 \(l\perp\) 平面 \(ADD_{1}A_{1}\)；
\item 设（1）中的直线 \(l\) 交 \(AC\) 于点 \(Q\). 求三棱锥 \(A_{1}-QC_{1}D\) 的体积. （锥体体积公式：\(V=\frac13Sh\)，其中 \(S\) 为底面面积，\(h\) 为高）
\end{enumerate}
\bitmapfigure[max width=0.54\linewidth]{img/2013/sichuan_liberal/q19_fig1.png}
\end{problem}

\begin{solution}
\begin{enumerate}
\item 平面 \(ABC\) 与平面 \(A_{1}BC\) 的交线是 \(BC\). 因此在平面 \(ABC\) 内过点 \(P\) 作 \(l\parallel BC\)，则 \(l\parallel\) 平面 \(A_{1}BC\).

又因为 \(D\) 是 \(BC\) 的中点，且 \(AB=AC\)，所以 \(AD\perp BC\). 由 \(l\parallel BC\)，得 \(l\perp AD\). 过 \(P\) 作直线 \(r\parallel AA_{1}\)，则 \(r\) 在平面 \(ADD_{1}A_{1}\) 内，且 \(r\) 与 \(AD\) 相交于 \(P\). 又因 \(AA_{1}\perp\) 平面 \(ABC\)，而 \(l\) 在平面 \(ABC\) 内，所以 \(r\perp l\). 因此 \(l\) 垂直于平面 \(ADD_{1}A_{1}\) 内过 \(P\) 的两条相交直线 \(AD\) 和 \(r\)，从而 \(l\perp\) 平面 \(ADD_{1}A_{1}\).

\item 设 \(\dfrac{AP}{AD}=t\)，其中 \(0<t<1\). 在三角形 \(ADC\) 中，\(P\in AD\)，且过 \(P\) 的直线 \(l\parallel BC\parallel DC\)，所以由相似三角形得 \(\dfrac{AQ}{AC}=\dfrac{AP}{AD}=t\).

取空间直角坐标系，使 \(A=(0,0,0)\)，\(B=(2,0,0)\)，\(C=(-1,\sqrt3,0)\)，并令 \(AA_{1}\) 方向为正方向，则 \(A_{1}=(0,0,1)\)，\(C_{1}=(-1,\sqrt3,1)\)，\(D=\left(\frac12,\frac{\sqrt3}{2},0\right)\)，\(Q=(-t,t\sqrt3,0)\).

以 \(A_{1}\) 为起点，三棱锥的三个棱向量为
\(\overrightarrow{A_{1}Q}=(-t,t\sqrt3,-1)\)，\(\overrightarrow{A_{1}C_{1}}=(-1,\sqrt3,0)\)，\(\overrightarrow{A_{1}D}=\left(\frac12,\frac{\sqrt3}{2},-1\right)\). 因此
\[
6V=\left|\det\begin{pmatrix}
-t&t\sqrt3&-1\\
-1&\sqrt3&0\\
\frac12&\frac{\sqrt3}{2}&-1
\end{pmatrix}\right|
=\left|(-t)(-\sqrt3)-t\sqrt3+(-1)(-\sqrt3)\right|=\sqrt3
\]
所以三棱锥 \(A_{1}-QC_{1}D\) 的体积为 \(V=\frac{\sqrt3}{6}\).
\end{enumerate}
\end{solution}

\begin{problem}
已知圆 \(C\) 的方程为 \(x^{2}+\left(y-4\right)^{2}=4\)，点 \(O\) 是坐标原点，直线 \(l:y=kx\) 与圆 \(C\) 交于 \(M\)，\(N\) 两点.
\begin{enumerate}
\item 求 \(k\) 的取值范围；
\item 设 \(Q(m,n)\) 是线段 \(MN\) 上的点，且 \(\dfrac{2}{|OQ|^{2}}=\dfrac{1}{|OM|^{2}}+\dfrac{1}{|ON|^{2}}\)，请将 \(n\) 表示为 \(m\) 的函数.
\end{enumerate}
\end{problem}

\begin{solution}
\begin{enumerate}
\item 将 \(y=kx\) 代入圆的方程，得
\((1+k^2)x^2-8kx+12=0\). 要使直线与圆有两个不同的交点，关于 \(x\) 的二次方程必须有两个不相等的实根，所以
\[\Delta=(-8k)^2-4(1+k^2)\cdot12=16(k^2-3)>0\]
因此 \(k\in(-\infty,-\sqrt3)\cup(\sqrt3,+\infty)\).

\item 设 \(M\)，\(N\) 的横坐标分别为 \(x_1\)，\(x_2\)，则 \(M=(x_1,kx_1)\)，\(N=(x_2,kx_2)\). 由韦达定理，
\(x_1+x_2=\frac{8k}{1+k^2}\)，\(x_1x_2=\frac{12}{1+k^2}\). 由于 \(Q\) 在线段 \(MN\) 上且在直线 \(y=kx\) 上，可写为 \(Q=(m,km)\)，所以 \(n=km\).

又有 \(|OQ|^2=(1+k^2)m^2\)，\(|OM|^2=(1+k^2)x_1^2\)，\(|ON|^2=(1+k^2)x_2^2\). 代入题设并约去正因子 \(1+k^2\)，得
\[
\frac{2}{m^2}=\frac{1}{x_1^2}+\frac{1}{x_2^2}=\frac{(x_1+x_2)^2-2x_1x_2}{(x_1x_2)^2}=\frac{5k^2-3}{18}
\]
因此 \(k^2=\frac{3(m^2+12)}{5m^2}\)，从而
\(n^2=k^2m^2=\frac35(m^2+12)\). 由于圆心为 \((0,4)\)，半径为 \(2\)，圆及弦 \(MN\) 均位于 \(x\) 轴上方，所以 \(n>0\)，得到
\(n=\sqrt{\frac35(m^{2}+12)}\).

由题设分母不为零，知 \(m\neq0\). 再由 \(|k|>\sqrt3\) 及上式中的 \(k^2\) 可得 \(0<m^2<3\). 还需验证这些 \(m\) 确实给出线段 \(MN\) 上的点. 令
\(F_k(x)=(1+k^2)x^2-8kx+12\)，其两个根为 \(x_1\)，\(x_2\). 对 \(0<m^2<3\)，并取与 \(n>0\) 相容的 \(k\) 的符号，有

\[
F_k(m)=(1+k^2)m^2-8km+12
=\frac{8(m^2+12)}5-8\sqrt{\frac{3(m^2+12)}5}<0
\]

这里最后一个不等号等价于 \(m^2<3\). 因为 \(F_k\) 的二次项系数为正，\(F_k(m)<0\) 说明 \(m\) 严格位于两根 \(x_1\)，\(x_2\) 之间，故 \(Q\) 确在线段 \(MN\) 内. 因此定义域为
\(-\sqrt3<m<0\) 或 \(0<m<\sqrt3\)，所求函数为
\(n=\sqrt{\frac35(m^2+12)}\).
\end{enumerate}
\end{solution}

\begin{problem}
已知函数 \(f(x)=\begin{cases}
x^2+2x+a,&x<0,\\
\ln x,&x>0,
\end{cases}\) 其中 \(a\) 是实数. 设 \(A(x_1,f(x_1))\)，\(B(x_2,f(x_2))\) 为该函数图象上的两点，且 \(x_1<x_2\).
\begin{enumerate}
\item 指出函数 \(f(x)\) 的单调区间；
\item 若函数 \(f(x)\) 的图象在点 \(A\)，\(B\) 处的切线互相垂直，且 \(x_2<0\)，证明：\(x_2-x_1\geqslant1\)；
\item 若函数 \(f(x)\) 的图象在点 \(A\)，\(B\) 处的切线重合，求 \(a\) 的取值范围.
\end{enumerate}
\end{problem}

\begin{solution}
\begin{enumerate}
\item 当 \(x<0\) 时，\(f'(x)=2x+2\)，所以函数在 \((-\infty,-1)\) 上单调递减，在 \((-1,0)\) 上单调递增；当 \(x>0\) 时，\(f'(x)=\frac1x>0\)，所以函数在 \((0,+\infty)\) 上单调递增.

\item 由 \(x_2<0\) 及 \(x_1<x_2\)，两点都在抛物线段上. 两处切线斜率分别为 \(2x_1+2\) 和 \(2x_2+2\)，垂直条件给出
\((2x_1+2)(2x_2+2)=-1\). 令 \(u=x_1+1\)，\(v=x_2+1\)，则 \(u<v<1\)，且 \(uv=-\frac14\). 因此 \(u<0<v\)，于是
\(x_2-x_1=v-u=v+(-u)\geqslant2\sqrt{v(-u)}=1\).

\item 若 \(x_1\)，\(x_2<0\)，两条切线斜率相等将导致 \(x_1=x_2\)，与 \(x_1<x_2\) 矛盾；若 \(x_1\)，\(x_2>0\)，同理由 \(\frac1{x_1}=\frac1{x_2}\) 也矛盾. 因此只能是 \(x_1<0<x_2\).

令 \(t=x_2>0\). 左侧抛物线在 \(x_1\) 处的切线为
\(y=(2x_1+2)x+a-x_1^2\)，右侧对数函数在 \(t\) 处的切线为 \(y=\frac1t x+\ln t-1\). 两切线重合，故斜率和截距分别相等：
\(2x_1+2=\frac1t\)，\(a-x_1^2=\ln t-1\). 从第一式得 \(x_1=\frac1{2t}-1\)，而 \(x_1<0\) 等价于 \(t>\frac12\). 代入第二式，得到
\[a=\left(\frac1{2t}-1\right)^2+\ln t-1=\frac1{4t^2}-\frac1t+\ln t\]
令 \(g(t)=\frac1{4t^2}-\frac1t+\ln t\)，则在 \(t>\frac12\) 时
\(g'(t)=\frac{2t^2+2t-1}{2t^3}>0\). 所以 \(g\) 在 \(\left(\frac12,+\infty\right)\) 上连续递增，且
\(\lim_{t\to\frac12+}g(t)=-1-\ln2\)，\(\lim_{t\to+\infty}g(t)=+\infty\). 因此
\(g\left(\frac12,+\infty\right)=(-1-\ln2,+\infty)\)，\(a\) 的取值范围为 \((-1-\ln2,+\infty)\).
\end{enumerate}
\end{solution}
